\begin{theorem}[ResidueMult]
\label{ResidueMult}
\lean{ResidueMult}
% at /home/runner/work/trough/trough/library/PrimeNumberTheoremAnd/ZetaBounds.lean:499.0-572.11
\leanok
If $f$ has a simple pole at $p$ with residue $A$, and $g$ is holomorphic near $p$, then the
  residue of $f \cdot g$ at $p$ is $A \cdot g(p)$. That is, we assume that
  $$
  f(s) = \frac{A}{s - p} + O(1)$$
  near $p$, and that $g$ is holomorphic near $p$. Then
  $$
  f(s) \cdot g(s) = \frac{A \cdot g(p)}{s - p} + O(1).$$
\end{theorem}
\begin{proof}
\leanok
Elementary calculation.
  $$
  f(s) * g(s) - \frac{A * g(p)}{s - p} =
  \left(f(s) * g(s) - \frac{A * g(s)}{s - p}\right)
  + \left(\frac{A * g(s) - A * g(p)}{s - p}\right).
  $$
  The first term is $g(s)(f(s) - \frac{A}{s - p})$, which is bounded near $p$ by the assumption
  on $f$
   and the fact that $g$ is holomorphic near $p$.
  The second term is $A$ times the log derivative of $g$ at $p$, which is bounded by the assumption
  that  $g$ is holomorphic.
\end{proof}
